\chapter{Expansión en series para las ecuaciones de las superficies cartesianas y demás cantidades de interés relacionadas}
\label{apx:ovoids_series}
\label{sec:series_expansion_cartesian_surfaces}

Con el propósito de estudiar el comportamiento local del sistema óptico en las cercanías del eje óptico, resulta natural desarrollar en series de potencias las ecuaciones paramétricas de la superficie cartesiana y las magnitudes geométricas y electromagnéticas derivadas de ella. Este procedimiento permite identificar la estructura paraxial del sistema, así como las primeras correcciones no paraxiales responsables de los efectos vectoriales que posteriormente se manifiestan en la formación de imagen.

Recordemos que la superficie dióptrica $\Sigma$ se describe mediante la parametrización


\begin{align}
z(\rho)
&=
\frac{(O+T\rho^2)\rho^2}
{1+S\rho^2+\sqrt{1+(2S-O^2G)\rho^2}},
\label{eq:cartesian_surface_parametrization_z}
\\[6pt]
r(\rho)
&=
\sqrt{\rho^2-z^2(\rho)}.
\label{eq:cartesian_surface_parametrization_r}
\end{align}

Los parámetros $G$, $O$, $T$ y $S$ que aparecen en \eqref{eq:cartesian_surface_parametrization_z} contienen la información geométrica del problema y se expresan en términos de los índices de refracción y de las posiciones axiales de los puntos conjugados $A$ y $A'$. Explícitamente,

\begin{align}
G
&=
\frac{(n_i^2z_i-n_0^2z_0)^2}
{n_in_0(n_iz_i-n_0z_0)(n_iz_0-n_0z_i)},
\label{eq:cartesian_parameter_G}
\\[6pt]
O
&=
\frac{n_iz_0-n_0z_i}
{z_iz_0(n_i-n_0)},
\label{eq:cartesian_parameter_O}
\\[6pt]
T
&=
\frac{(n_i-n_0)(n_i+n_0)^2}
{4n_in_0z_iz_0(n_iz_i-n_0z_0)},
\label{eq:cartesian_parameter_T}
\\[6pt]
S
&=
\frac{(n_i+n_0)(n_i^2z_i-n_0^2z_0)}
{2n_in_0z_iz_0(n_iz_i-n_0z_0)}.
\label{eq:cartesian_parameter_S}
\end{align}

A partir de \eqref{eq:cartesian_surface_parametrization_z}, expandiendo alrededor de $\rho=0$, se obtiene

\begin{align}
z(\rho)
=
{}&
\frac{O}{2}\rho^2
+
\left(
\frac{T}{2}
-
\frac{OS}{2}
+
\frac{O^3G}{8}
\right)\rho^4
\nonumber
\\[6pt]
&
+
\left(
\frac{OS^2}{2}
-
\frac{O^3GS}{4}
+
\frac{O^5G^2}{16}
-
\frac{TS}{2}
+
\frac{TO^2G}{8}
\right)\rho^6
+
O(\rho^8).
\label{eq:z_rho_series_expansion}
\end{align}

Luego, sustituyendo \eqref{eq:z_rho_series_expansion} en la relación radial \eqref{eq:cartesian_surface_parametrization_r}, se obtiene

\begin{align}
r(\rho)
=
{}&
\rho
-
\frac{O^2}{8}\rho^3
\nonumber
\\[6pt]
&
+
\left(
-\frac{OT}{4}
+
\frac{O^2S}{4}
-
\frac{O^4G}{16}
-
\frac{O^4}{128}
\right)\rho^5
+
O(\rho^7).
\label{eq:r_rho_series_expansion}
\end{align}


O remplazando los parametros G, O, T, S. 

\begin{align}
z(\rho)
&=
\frac{n_i z_0 - n_0 z_i}
{2 z_i z_0 (n_i-n_0)}
\rho^2
+
\frac{
n_i n_0 (z_0-z_i)^2 (n_i z_i - n_0 z_0)
}{
8 z_i^3 z_0^3 (n_i-n_0)^3
}
\rho^4
\\[6pt]
&\quad
+
\frac{
n_i^2 z_i - n_0^2 z_0
}{
32 n_i^2 n_0^2 z_i^5 z_0^5
(n_i-n_0)^5
(n_i z_i - n_0 z_0)^2
}
\\[4pt]
&\qquad\times
\Big[
2
(n_i z_0 - n_0 z_i)^3
(n_i^2 z_i - n_0^2 z_0)^3
\\
&\qquad\qquad
-
4
(n_i z_0 - n_0 z_i)^2
(n_i^2 z_i - n_0^2 z_0)^2
(n_i-n_0)^2
z_i z_0
(n_i+n_0)
\\
&\qquad\qquad
+
5
(n_i z_0 - n_0 z_i)
(n_i^2 z_i - n_0^2 z_0)
(n_i-n_0)^4
z_i^2 z_0^2
(n_i+n_0)^2
\\
&\qquad\qquad
-
2
(n_i-n_0)^6
z_i^3 z_0^3
(n_i+n_0)^3
\Big]\rho^6
+
O(\rho^8)
\end{align}


\begin{align}
r(\rho)
&=
\rho
-
\frac{(n_i z_0-n_0 z_i)^2}
{8z_i^2z_0^2(n_i-n_0)^2}
\rho^3
\\[6pt]
&\quad
+
\frac{(n_i z_0-n_0 z_i)}
{128 z_0^4 z_i^4 (n_i-n_0)^4}
\Big[
n_0^3z_i^3
+8n_0^2n_i z_0^3
-16n_0^2n_i z_0^2z_i
\\
&\qquad\qquad
+5n_0^2n_i z_0z_i^2
-5n_0n_i^2 z_0^2z_i
+16n_0n_i^2 z_0z_i^2
-8n_0n_i^2 z_i^3
-n_i^3z_0^3
\Big]\rho^5
+
O(\rho^7)
\end{align}




El término dominante de \eqref{eq:z_rho_series_expansion} es

\begin{equation}
z(\rho)\sim \frac{O}{2}\rho^2.
\label{eq:vertex_paraxial_curvature}
\end{equation}

Por tanto, en virtud de \eqref{eq:vertex_paraxial_curvature}, el parámetro $O$ controla la curvatura en el vértice de la superficie. En efecto, como \eqref{eq:r_rho_series_expansion} implica que $r(\rho)\sim \rho$, se tiene localmente

\begin{equation}
z(r)\sim \frac{O}{2}r^2,
\label{eq:vertex_curvature_sag_form}
\end{equation}

lo cual identifica a $O$ como el parámetro de curvatura paraxial de la superficie.

Sea ahora

\begin{equation}
Q(\rho,\phi)
=
r(\rho)\,\uvec{\rho}
+
z(\rho)\,\uvec{z}
\in\Sigma
\label{eq:surface_point_Q}
\end{equation}

un punto sobre la superficie cartesiana, y considérense los puntos conjugados

\begin{equation}
A=(0,0,z_0),
\qquad
A'=(0,0,z_i).
\label{eq:conjugate_points_A_Aprime}
\end{equation}

Las distancias geométricas desde los puntos conjugados \eqref{eq:conjugate_points_A_Aprime} hasta el punto superficial \eqref{eq:surface_point_Q} se definen como

\begin{equation}
\ell_0=AQ,
\qquad
\ell_i=A'Q.
\label{eq:geometrical_distances_l0_li}
\end{equation}

Usando la identidad exacta que se deduce de \eqref{eq:cartesian_surface_parametrization_r},

\begin{equation}
r^2(\rho)+z^2(\rho)=\rho^2,
\label{eq:identity_r2_z2_rho2}
\end{equation}

las distancias \eqref{eq:geometrical_distances_l0_li} satisfacen

\begin{align}
\ell_0^2
&=
r^2(\rho)+\bigl(z(\rho)-z_0\bigr)^2
=
\rho^2+z_0^2-2z_0z(\rho),
\label{eq:l0_squared_exact}
\\[6pt]
\ell_i^2
&=
r^2(\rho)+\bigl(z(\rho)-z_i\bigr)^2
=
\rho^2+z_i^2-2z_iz(\rho).
\label{eq:li_squared_exact}
\end{align}

Sustituyendo \eqref{eq:z_rho_series_expansion} en \eqref{eq:l0_squared_exact} y \eqref{eq:li_squared_exact}, y reemplazando completamente los parámetros $G$, $O$, $T$ y $S$ mediante \eqref{eq:cartesian_parameter_G}--\eqref{eq:cartesian_parameter_S}, se obtiene, bajo la convención axial $z_0<0$ y $z_i>0$,

\begin{align}
\ell_0(\rho)
=
{}&
-z_0
-
\frac{n_i(z_0-z_i)}
{2z_0z_i(n_0-n_i)}
\rho^2
\nonumber
\\[4pt]
&
+
\frac{
n_i(z_0-z_i)^2
\left(
n_0^2z_0-n_i^2z_i
\right)
}
{
8z_0^3z_i^3(n_0-n_i)^3
}
\rho^4
\nonumber
\\[4pt]
&
-
\frac{
n_i(z_0-z_i)^3
\left(
n_0^2z_0-n_i^2z_i
\right)^2
}
{
16z_0^5z_i^5(n_0-n_i)^5
}
\rho^6
\nonumber
\\[4pt]
&
+
\frac{
5n_i(z_0-z_i)^4
\left(
n_0^2z_0-n_i^2z_i
\right)^3
}
{
128z_0^7z_i^7(n_0-n_i)^7
}
\rho^8
+
O(\rho^{10}),
\label{eq:l0_series_expansion}
\end{align}

y

\begin{align}
\ell_i(\rho)
=
{}&
z_i
+
\frac{n_0(z_0-z_i)}
{2z_0z_i(n_0-n_i)}
\rho^2
\nonumber
\\[4pt]
&
-
\frac{
n_0(z_0-z_i)^2
\left(
n_0^2z_0-n_i^2z_i
\right)
}
{
8z_0^3z_i^3(n_0-n_i)^3
}
\rho^4
\nonumber
\\[4pt]
&
+
\frac{
n_0(z_0-z_i)^3
\left(
n_0^2z_0-n_i^2z_i
\right)^2
}
{
16z_0^5z_i^5(n_0-n_i)^5
}
\rho^6
\nonumber
\\[4pt]
&
-
\frac{
5n_0(z_0-z_i)^4
\left(
n_0^2z_0-n_i^2z_i
\right)^3
}
{
128z_0^7z_i^7(n_0-n_i)^7
}
\rho^8
+
O(\rho^{10}).
\label{eq:li_series_expansion}
\end{align}

Las expansiones \eqref{eq:l0_series_expansion} y \eqref{eq:li_series_expansion} verifican término a término la condición cartesiana de camino óptico constante,

\begin{equation}
n_0\ell_0+n_i\ell_i
=
\text{constante}.
\label{eq:cartesian_optical_path_condition}
\end{equation}

En efecto, al sustituir \eqref{eq:l0_series_expansion} y \eqref{eq:li_series_expansion} en \eqref{eq:cartesian_optical_path_condition}, los términos de orden $\rho^2$, $\rho^4$, $\rho^6$ y $\rho^8$ se cancelan, de modo que

\begin{equation}
n_0\ell_0+n_i\ell_i
=
-n_0z_0+n_iz_i
+
O(\rho^{10}).
\label{eq:optical_path_constant_local}
\end{equation}

La cancelación expresada en \eqref{eq:optical_path_constant_local} es la manifestación local del estigmatismo impuesto por la superficie cartesiana.

A partir de las cantidades geométricas anteriores se obtienen los cosenos de incidencia y transmisión que intervienen en los coeficientes de Fresnel. Con la convención orientada adoptada,

\begin{equation}
\cos\theta_i(0)=-1,
\qquad
\cos\theta_t(0)=1.
\label{eq:cos_theta_axis_limit}
\end{equation}

Expandiendo los cosenos alrededor de $\rho=0$, y usando las expansiones de $\ell_0$ y $\ell_i$ dadas en \eqref{eq:l0_series_expansion} y \eqref{eq:li_series_expansion}, se encuentra

\begin{align}
\cos\theta_i
=
{}&
-1
+
\frac{n_i^2(z_0-z_i)^2}
{2z_0^2z_i^2(n_0-n_i)^2}
\rho^2
\nonumber
\\[4pt]
&
-
\frac{
n_i^2(z_0-z_i)^2
\left(
4n_0^2z_0^2
-4n_0^2z_0z_i
+n_0^2z_i^2
-2n_0n_iz_0z_i
-2n_i^2z_0z_i
+3n_i^2z_i^2
\right)
}
{8z_0^4z_i^4(n_0-n_i)^4}
\rho^4
+
O(\rho^6),
\label{eq:cos_theta_i_series}
\end{align}

mientras que

\begin{align}
\cos\theta_t
=
{}&
1
-
\frac{n_0^2(z_0-z_i)^2}
{2z_0^2z_i^2(n_0-n_i)^2}
\rho^2
\nonumber
\\[4pt]
&
+
\frac{
n_0^2(z_0-z_i)^2
\left(
3n_0^2z_0^2
-2n_0^2z_0z_i
-2n_0n_iz_0z_i
+n_i^2z_0^2
-4n_i^2z_0z_i
+4n_i^2z_i^2
\right)
}
{8z_0^4z_i^4(n_0-n_i)^4}
\rho^4
+
O(\rho^6).
\label{eq:cos_theta_t_series}
\end{align}

Los signos dominantes en \eqref{eq:cos_theta_i_series} y \eqref{eq:cos_theta_t_series} son consistentes con \eqref{eq:cos_theta_axis_limit}. La diferencia de signo corresponde a la orientación de la normal respecto a los rayos incidente y transmitido.

Los coeficientes de Fresnel pueden escribirse en términos de las cantidades geométricas ya desarrolladas. Para ello definimos

\begin{equation}
\eta=\frac{n_i}{n_0}.
\label{eq:eta_definition}
\end{equation}

Con esta notación, los coeficientes exactos son

\begin{align}
t_s
&=
\frac{
2\left[
z'\left(\rho^2-zz_0\right)-\rho\left(z-z_0\right)
\right]
}
{
\left[
z'\left(\rho^2-zz_0\right)-\rho\left(z-z_0\right)
\right]
+
\eta
\frac{\ell_0}{\ell_i}
\left[
z'\left(\rho^2-zz_i\right)-\rho\left(z-z_i\right)
\right]
},
\label{eq:fresnel_ts_exact}
\\[8pt]
t_p
&=
\frac{
2\left[
z'\left(\rho^2-zz_0\right)-\rho\left(z-z_0\right)
\right]
}
{
\eta
\left[
z'\left(\rho^2-zz_0\right)-\rho\left(z-z_0\right)
\right]
+
\frac{\ell_0}{\ell_i}
\left[
z'\left(\rho^2-zz_i\right)-\rho\left(z-z_i\right)
\right]
}.
\label{eq:fresnel_tp_exact}
\end{align}

En \eqref{eq:fresnel_ts_exact} y \eqref{eq:fresnel_tp_exact}, tanto $z(\rho)$ como $\ell_0(\rho)$ y $\ell_i(\rho)$ deben entenderse mediante sus desarrollos respectivos \eqref{eq:z_rho_series_expansion}, \eqref{eq:l0_series_expansion} y \eqref{eq:li_series_expansion}. Al expandir \eqref{eq:fresnel_ts_exact} y \eqref{eq:fresnel_tp_exact}, se obtiene

\begin{align}
t_s(\rho)
=
{}&
\frac{2n_0}{n_0-n_i}
-
\frac{
n_0n_i(n_0+n_i)(z_0-z_i)^2
}
{
z_0^2z_i^2(n_0-n_i)^3
}
\rho^2
\nonumber
\\[4pt]
&
+
\frac{
n_0n_i(n_0+n_i)(z_0-z_i)^2
(n_0z_0-n_iz_i)
(3n_0z_0-2n_0z_i+2n_iz_0-3n_iz_i)
}
{
4z_0^4z_i^4(n_0-n_i)^5
}
\rho^4
+
O(\rho^6),
\label{eq:fresnel_ts_series}
\end{align}

y

\begin{align}
t_p(\rho)
=
{}&
\frac{2n_0}{n_i-n_0}
-
\frac{
n_0^2(n_0+n_i)(z_0-z_i)^2
}
{
z_0^2z_i^2(n_0-n_i)^3
}
\rho^2
\nonumber
\\[4pt]
&
+
\frac{
n_0^2(n_0+n_i)(z_0-z_i)^2
}
{
4z_0^4z_i^4(n_0-n_i)^5
}
\nonumber
\\
&\quad\times
\Big[
n_0^2z_0^2
+2n_0^2z_0z_i
-2n_0^2z_i^2
-2n_0n_iz_0^2
+2n_0n_iz_0z_i
-2n_0n_iz_i^2
\nonumber
\\
&\qquad
-2n_i^2z_0^2
+2n_i^2z_0z_i
+n_i^2z_i^2
\Big]\rho^4
+
O(\rho^6).
\label{eq:fresnel_tp_series}
\end{align}

De las expansiones \eqref{eq:fresnel_ts_series} y \eqref{eq:fresnel_tp_series} se deduce que las combinaciones que aparecen naturalmente en la base circular satisfacen

\begin{align}
t_p+t_s
&=
-
\frac{
n_0(n_0+n_i)^2(z_0-z_i)^2
}
{
z_0^2z_i^2(n_0-n_i)^3
}
\rho^2
+
O(\rho^4),
\label{eq:tp_plus_ts_series_paraxial}
\\[8pt]
t_p-t_s
&=
\frac{4n_0}{n_i-n_0}
-
\frac{
n_0(n_0+n_i)(z_0-z_i)^2
}
{
z_0^2z_i^2(n_0-n_i)^2
}
\rho^2
+
O(\rho^4).
\label{eq:tp_minus_ts_paraxial}
\end{align}



Por tanto, si

\begin{align}
\mathscr F\{T_L\}
&=
\mathscr F
\left\{
(t_p+t_s)E_L
+
(t_p-t_s)e^{2i\phi'}E_R
\right\},
\label{eq:FT_TL_exact_circular}
\\[6pt]
\mathscr F\{T_R\}
&=
\mathscr F
\left\{
(t_p-t_s)e^{-2i\phi'}E_L
+
(t_p+t_s)E_R
\right\},
\label{eq:FT_TR_exact_circular}
\end{align}

entonces, sustituyendo \eqref{eq:tp_plus_ts_series_paraxial} y \eqref{eq:tp_minus_ts_paraxial} en \eqref{eq:FT_TL_exact_circular} y \eqref{eq:FT_TR_exact_circular}, y usando que al orden retenido $\rho'^2=r'^2+O(r'^4)$, se obtiene

\begin{align}
\mathscr F\{T_L\}
\simeq{}&
-
2\pi
\frac{
n_0(n_0+n_i)^2(z_0-z_i)^2
}
{
z_0^2z_i^2(n_0-n_i)^3
}
\mathscr H_0[r'^2E_L]
\\[6pt]
&-
\frac{8\pi n_0}{n_i-n_0}
e^{2i\varphi}
\mathscr H_2[E_R]
\\[6pt]
&+
2\pi
\frac{
n_0(n_0+n_i)(z_0-z_i)^2
}
{
z_0^2z_i^2(n_0-n_i)^2
}
e^{2i\varphi}
\mathscr H_2[r'^2E_R],
\label{eq:FT_TL_paraxial}
\\[10pt]
\mathscr F\{T_R\}
\simeq{}&
-
\frac{8\pi n_0}{n_i-n_0}
e^{-2i\varphi}
\mathscr H_2[E_L]
\\[6pt]
&+
2\pi
\frac{
n_0(n_0+n_i)(z_0-z_i)^2
}
{
z_0^2z_i^2(n_0-n_i)^2
}
e^{-2i\varphi}
\mathscr H_2[r'^2E_L]
\\[6pt]
&-
2\pi
\frac{
n_0(n_0+n_i)^2(z_0-z_i)^2
}
{
z_0^2z_i^2(n_0-n_i)^3
}
\mathscr H_0[r'^2E_R].
\label{eq:FT_TR_paraxial}
\end{align}

La estructura de \eqref{eq:FT_TL_paraxial} y \eqref{eq:FT_TR_paraxial} exhibe la jerarquía local del acoplamiento vectorial: los términos diagonales aparecen en el siguiente orden paraxial, mientras que los términos de conversión circular permanecen en el orden dominante. Esta separación será esencial para identificar, en el campo focal, cuáles contribuciones pertenecen al límite escalar-paraxial y cuáles corresponden a correcciones vectoriales inducidas por la geometría local de la superficie cartesiana.
